
The routines in this section provide initialization and memory management for the
\ac{GPU} symmetric heap. Runtime capability reporting is handled by
\FUNC{shmem\_init\_thread}.

\textbf{Prerequisite:} Except for \FUNC{shmem\_init\_thread}, all routines in
this and subsequent sections require that the \openshmem library has been
initialized with GPU support by passing \CONST{SHMEM\_DEVICE\_HOST\_INIT},
\CONST{SHMEM\_DEVICE\_KERNEL\_INIT}, or both to \FUNC{shmem\_init\_thread}, and
that the corresponding capability was returned in \VAR{provided}. Calling any
\texttt{shmemg\_} routine, or using a routine defined in \openshmem 1.6 with
GPU-attached memory operands, without the required capability having been
returned in \VAR{provided} results in \textbf{undefined behavior}.

The capability required by each group of routines is summarized below.

\begin{longtable}{|p{7.5cm}|p{7cm}|}
\hline
\textbf{Routine group} & \textbf{Required capability in \VAR{provided}} \\
\hline
\endfirsthead
\hline
\textbf{Routine group} & \textbf{Required capability in \VAR{provided}} \\
\hline
\endhead
\ac{GPU} symmetric heap management (\cref{sec:init_memmgmt}) &
\CONST{SHMEM\_DEVICE\_HOST\_INIT} or \CONST{SHMEM\_DEVICE\_KERNEL\_INIT} \\
\hline
GPU-aware communication (\cref{sec:gpu_aware}) &
\CONST{SHMEM\_DEVICE\_HOST\_INIT} \\
\hline
GPU-centric communication (\cref{sec:gpu_centric}) &
\CONST{SHMEM\_DEVICE\_KERNEL\_INIT} \\
\hline
\end{longtable}

\subsection{SHMEM\_INIT\_THREAD}\label{subsec:shmem_init_thread}

Initializes the \openshmem library with support for the requested thread level
and \ac{GPU} communication capabilities.

\begin{apidefinition}
\begin{Csynopsis}
int shmem_init_thread(int requested, int *provided);
\end{Csynopsis}

\begin{apiarguments}
\apiargument{IN}{requested}{The thread level and GPU support requested by the
application. The value is a bitwise OR of a thread level constant and zero or
more GPU support constants.}
\apiargument{OUT}{provided}{The thread level and GPU support provided by the
implementation.}
\end{apiarguments}

\apidescription{
The thread level portion of \VAR{requested} and \VAR{provided} shall be one of
\CONST{SHMEM\_THREAD\_SINGLE}, \CONST{SHMEM\_THREAD\_FUNNELED},
\CONST{SHMEM\_THREAD\_SERIALIZED}, or \CONST{SHMEM\_THREAD\_MULTIPLE}, as
defined in \openshmem 1.6.

The GPU support portion of \VAR{requested} and \VAR{provided} may include one or
both of the following constants, combined with the thread level using bitwise OR:
\begin{itemize}
\item \CONST{SHMEM\_DEVICE\_HOST\_INIT} --- requests GPU-aware communication support
\item \CONST{SHMEM\_DEVICE\_KERNEL\_INIT} --- requests GPU-centric communication support
\end{itemize}

\FUNC{shmem\_init\_thread} is extended to accept GPU support options in the
\VAR{requested} argument. When any GPU support option is included in
\VAR{requested}, the implementation shall establish the PE-to-GPU association and
initialize the \ac{GPU} symmetric heap. When \CONST{SHMEM\_DEVICE\_HOST\_INIT}
is included in \VAR{requested}, the implementation shall prepare the runtime for
GPU-aware communication. When \CONST{SHMEM\_DEVICE\_KERNEL\_INIT} is included in
\VAR{requested}, the implementation shall prepare the runtime for GPU-centric
communication operations invoked from within GPU kernels.

The \VAR{provided} argument returns the thread level and GPU support actually
provided by the implementation. The GPU support bits in \VAR{provided} shall
reflect only the capabilities that the implementation has successfully
initialized. An implementation may provide a subset of the requested GPU support.

If no GPU support constants are included in \VAR{requested}, the behavior of
\FUNC{shmem\_init\_thread} is unchanged from \openshmem 1.6.

The application shall select the \ac{GPU} for association with the \ac{PE} prior
to calling \FUNC{shmem\_init\_thread}. The implementation shall use the \ac{GPU}
that has been made current by the execution environment at the time of the call.

All \acp{PE} shall provide the same GPU support constants in \VAR{requested}. If
differing values are used, the behavior is undefined.
}

\apireturnvalues{
\FUNC{shmem\_init\_thread} returns 0 upon success; otherwise, it returns a
nonzero value.
}

\apinotes{
The application should compare the GPU support bits in \VAR{provided} against the
requested bits to determine whether the desired GPU capabilities are available:
\begin{itemize}
\item If \CONST{SHMEM\_DEVICE\_HOST\_INIT} is \textbf{not} set in
  \VAR{provided} --- do not use GPU-attached memory buffers in \openshmem
  operations
\item If \CONST{SHMEM\_DEVICE\_KERNEL\_INIT} is \textbf{not} set in
  \VAR{provided} --- do not invoke \openshmem routines from GPU kernels
\end{itemize}
}
\end{apidefinition}

\subsection{Memory Management}\label{subsec:gpu_memmgmt}

\subsubsection{SHMEMG\_MALLOC}\label{subsec:shmemg_malloc}

Collectively allocate a block of memory from the \ac{GPU} symmetric heap.

\begin{apidefinition}
\begin{Csynopsis}
void *shmemg_malloc(size_t size);
\end{Csynopsis}

\begin{apiarguments}
\apiargument{IN}{size}{The size, in bytes, of a block to be allocated from the
\ac{GPU} symmetric heap.}
\end{apiarguments}

\apidescription{
The \FUNC{shmemg\_malloc} routine is a collective operation on the world team
that returns the symmetric address of a block of at least \VAR{size} bytes from
the \ac{GPU} symmetric heap. The returned address shall be suitably aligned so
that it may be assigned to a pointer to any type of object. The allocated memory
resides in GPU-attached memory on the \ac{GPU} associated with the calling
\ac{PE}.

When \VAR{size} is zero, the routine performs no action and returns a null
pointer. Otherwise, \FUNC{shmemg\_malloc} calls a procedure that is semantically
equivalent to \FUNC{shmem\_barrier\_all} on exit. This ensures that all \acp{PE}
participate in the memory allocation, and that the memory on other \acp{PE} can
be used as soon as the local \ac{PE} returns.

The value of the \VAR{size} argument must be identical on all \acp{PE};
otherwise, the behavior is undefined.
}

\apireturnvalues{
\FUNC{shmemg\_malloc} returns the symmetric address of the allocated space;
otherwise, it returns a null pointer.
}
\end{apidefinition}

\subsubsection{SHMEMG\_CALLOC}\label{subsec:shmemg_calloc}

Collectively allocate a zeroed block of memory from the \ac{GPU} symmetric heap.

\begin{apidefinition}
\begin{Csynopsis}
void *shmemg_calloc(size_t count, size_t size);
\end{Csynopsis}

\begin{apiarguments}
\apiargument{IN}{count}{The number of elements to allocate.}
\apiargument{IN}{size}{The size, in bytes, of each element to allocate.}
\end{apiarguments}

\apidescription{
The \FUNC{shmemg\_calloc} routine is a collective operation on the world team
that allocates a region of remotely accessible memory from the \ac{GPU} symmetric
heap for an array of \VAR{count} objects of \VAR{size} bytes each. The returned
pointer refers to the lowest byte address of the allocated memory. The space is
initialized to all bits zero. The allocated memory resides in GPU-attached memory
on the \ac{GPU} associated with the calling \ac{PE}.

If the allocation succeeds, the returned pointer shall be suitably aligned so
that it may be assigned to a pointer to any type of object. If the allocation
does not succeed, or either \VAR{count} or \VAR{size} is 0, the return value is
a null pointer.

The values of \VAR{count} and \VAR{size} shall each be equal across all
\acp{PE}; otherwise, the behavior is undefined.

When \VAR{count} or \VAR{size} is 0, the routine returns without performing a
barrier. Otherwise, \FUNC{shmemg\_calloc} calls a procedure that is semantically
equivalent to \FUNC{shmem\_barrier\_all} on exit.
}

\apireturnvalues{
\FUNC{shmemg\_calloc} returns a pointer to the lowest byte address of the
allocated space; otherwise, it returns a null pointer.
}
\end{apidefinition}

\subsubsection{SHMEMG\_ALIGN}\label{subsec:shmemg_align}

Collectively allocate a block of memory with a specified alignment from the
\ac{GPU} symmetric heap.

\begin{apidefinition}
\begin{Csynopsis}
void *shmemg_align(size_t alignment, size_t size);
\end{Csynopsis}

\begin{apiarguments}
\apiargument{IN}{alignment}{Byte alignment of the block allocated from the
\ac{GPU} symmetric heap.}
\apiargument{IN}{size}{The size, in bytes, of a block to be allocated from the
\ac{GPU} symmetric heap.}
\end{apiarguments}

\apidescription{
The \FUNC{shmemg\_align} routine is a collective operation on the world team that
allocates a block from the \ac{GPU} symmetric heap with a byte alignment
specified by the \VAR{alignment} argument. The value of \VAR{alignment} shall be
a multiple of \texttt{sizeof(void *)} that is also a power of two; otherwise, the
behavior is undefined. The allocated memory resides in GPU-attached memory on the
\ac{GPU} associated with the calling \ac{PE}.

When \VAR{size} is zero, the routine performs no action and returns a null
pointer. Otherwise, \FUNC{shmemg\_align} calls a procedure that is semantically
equivalent to \FUNC{shmem\_barrier\_all} on exit.

The values of \VAR{alignment} and \VAR{size} must be identical on all \acp{PE};
otherwise, the behavior is undefined.
}

\apireturnvalues{
\FUNC{shmemg\_align} returns an aligned symmetric address whose value is a
multiple of \VAR{alignment}; otherwise, it returns a null pointer.
}
\end{apidefinition}

\subsubsection{SHMEMG\_FREE}\label{subsec:shmemg_free}

Collectively deallocate memory from the \ac{GPU} symmetric heap.

\begin{apidefinition}
\begin{Csynopsis}
void shmemg_free(void *ptr);
\end{Csynopsis}

\begin{apiarguments}
\apiargument{IN}{ptr}{Symmetric address of an object in the \ac{GPU} symmetric
heap.}
\end{apiarguments}

\apidescription{
The \FUNC{shmemg\_free} routine is a collective operation on the world team that
deallocates the block to which \VAR{ptr} points, making it available for further
allocation. If \VAR{ptr} is a null pointer, no action is performed; otherwise,
\FUNC{shmemg\_free} calls a procedure that is semantically equivalent to
\FUNC{shmem\_barrier\_all} on entry.

It is the application's responsibility to ensure that no communication operations
involving the given memory block are pending on any communication context prior
to calling \FUNC{shmemg\_free}.

The value of \VAR{ptr} must be identical on all \acp{PE}; otherwise, the
behavior is undefined.
}

\apireturnvalues{
None.
}
\end{apidefinition}

\clearpage

